
\section{结论}
我们提出了一种新颖的框架，在提升 LLM 推理能力的同时实现参数压缩。这是通过利用关于 LLM 中多头注意力的领域知识与实证依据，结合独特的多头张量化与 Tucker 分解的一种特殊变体来实现的。借此，我们的框架探索了按照注意力头之间共享的高维低秩结构，对单个 Transformer 层中每个注意力头的权重进行结构性去噪。由此确保每个注意力头的权重在由公共 Tucker 因子矩阵刻画的同一高维子空间内编码不同的信息。我们的方法已被证明能够在仅编码器与仅解码器的 LLM（参数规模从数亿到数十亿）中实现参数压缩与推理增强。此外，我们的框架可与现有提升 LLM 推理能力的方法（如基于 FFN 权重去噪的方法）联合使用，以获得进一步的性能增益。通过消融研究，我们验证了所提多头张量化过程带来的优势与性能提升。
\paragraph{局限性}
与其他现有权重去噪方法类似，我们发现对于不同的数据集，本文方法在不同的超参数设置下取得最佳结果。我们未来的工作将致力于为本文方法及其他现有方法找到统一且可泛化的超参数设置。


\bibliography{reference}
\bibliographystyle{ieeetr}

\end{document}
